\begin{table}[H]
\centering
\footnotesize
\caption{Las cinco confusiones dirigidas más frecuentes de cada evaluación; los empates se ordenan por identificador de clase.}
\label{tab:confusions}
\begin{tabularx}{\linewidth}{l Y Y r}
\toprule
\textbf{Modelo / partición} & \textbf{Género real} & \textbf{Género predicho} & \textbf{Conteo}\\
\midrule
R / val & \textit{Amanita} & \textit{Lycoperdon} & 22\\
R / val & \textit{Ganoderma} & \textit{Trametes} & 14\\
R / val & \textit{Trametes} & \textit{Ganoderma} & 10\\
R / val & \textit{Trametes} & \textit{Lycoperdon} & 9\\
R / val & \textit{Ganoderma} & \textit{Xylaria} & 6\\
C / val & \textit{Trametes} & \textit{Ganoderma} & 10\\
C / val & \textit{Auricularia} & \textit{Ganoderma} & 7\\
C / val & \textit{Lycoperdon} & \textit{Amanita} & 7\\
C / val & \textit{Auricularia} & \textit{Trametes} & 6\\
C / val & \textit{Auricularia} & \textit{Pleurotus} & 5\\
R / test & \textit{Amanita} & \textit{Lycoperdon} & 26\\
R / test & \textit{Ganoderma} & \textit{Trametes} & 13\\
R / test & \textit{Trametes} & \textit{Ganoderma} & 11\\
R / test & \textit{Auricularia} & \textit{Xylaria} & 10\\
R / test & \textit{Trametes} & \textit{Lycoperdon} & 7\\
\bottomrule
\end{tabularx}
\source{reports/s6/comparison/confusion_pairs_top5.csv}
\end{table}
